% ======================================================================
%  Sections 4.6.3 and 4.6.4 --- Test configurations, and Metrics.
%  KEEP THE METRICS SUBSECTION. It needs three corrections, one of them
%  load-bearing for the whole chapter.
%
%  A. THE FAILURE DEFINITION DECLARES EVERY RUN A FAILURE.
%     "a run being declared to have failed if the stopping threshold is
%     not reached within the iteration cap." The threshold is
%     ||e||^2 < 10^4; 4.7.1 reports the feature error settling three to
%     four orders of magnitude above it; no run reaches it. This is the
%     consequent edit promised in the 4.6.2 correction, and it cannot be
%     left: as written, the protocol convicts its own results.
%
%  B. A METRIC THE CHAPTER NEVER REPORTS. "the number of iterations to
%     convergence" is listed, then excused as secondary, and then never
%     given anywhere in 4.7. This is the same defect as 4.3.2's promise to
%     characterise the admissible occlusion fraction and 4.5's claim to
%     monitor the conditioning -- both of which we removed. Removing this
%     one keeps the chapter consistent with itself.
%
%  C. THE FEATURE ERROR IS LISTED AS A CRITERION, AND THE CHAPTER'S
%     SHARPEST RESULT IS THAT IT IS NOT ONE.
%     Under occlusion Variant 1 attains the SMALLEST final feature error
%     of the three, 1.551e10 against 1.573e10 for Variant 3, while its
%     pose error is some four hundred times larger, 1291 um against
%     3.12 um. The feature error is exactly the quantity the occluder
%     biases, so a reader who takes this subsection at its word and reads
%     the ||e||^2 column of Table 4.x as a performance ranking will draw
%     the reverse of the chapter's conclusion. The metric has to be kept
%     -- it is the convergence evidence -- but its standing has to be
%     declared here, where the metrics are declared.
%
%  D. NO \label. The corrected 4.6.2 refers to Section~\ref{sec:metrics}.
%
%  E. 4.6.3 was an empty heading. Drafted below from the run
%     configuration. Check the three numbers marked  <-- CONFIRM.
% ======================================================================

\subsection{Test configurations}
\label{sec:test-configurations}

Two configurations are used, differing only in the target texture presented
to the camera during servoing.

In the \emph{nominal} configuration the reference image and the images
acquired during servoing are both rendered from the same unoccluded texture.
No measurement is corrupted, and the weighting has nothing to reject: the
configuration tests whether the robust schemes preserve the accuracy of the
baseline when robustness is not called for, which is the precondition for
using them at all.

In the \emph{occluded} configuration the reference image is rendered from the
unoccluded texture while the images acquired during servoing are rendered
from a copy of it carrying an opaque, untextured elliptical region. The
occluder is painted into the target texture rather than introduced as a
separate object, so it is planar with the target and its projection scales
with the camera depth: it covers $7.4\%$ of the measurements at the initial
pose and $5.80\%$ at the desired pose, the reduction following from the
retreat in depth described below. Since the reference is taken from the clean
texture, every measurement falling inside the occluder is in error, and the
error is systematic rather than random --- which is what distinguishes this
test from the addition of noise.

The poses are identical in both configurations. The camera starts at
\begin{equation*}
    \mathbf{t}_{0}=(40,\,-40,\,1100)\,\mathrm{mm},
    \qquad
    \boldsymbol{\omega}_{0}=(12,\,-12,\,6)^{\circ},
\end{equation*}
and the desired pose is fronto-parallel to the target at
\begin{equation*}
    \mathbf{t}^{*}=(0,\,0,\,1300)\,\mathrm{mm},
    \qquad
    \boldsymbol{\omega}^{*}=(0,\,0,\,0)^{\circ} .
\end{equation*}
The displacement combines a lateral translation, a retreat in depth and a
rotation about all three axes, and induces a mean image displacement of some
$41$ pixels, reaching $65$ pixels in the corners of the field of view.
%%% <-- CONFIRM these two pixel figures against your run log before
%%% submission; the poses and the two occlusion percentages are measured.
This is large enough that the linearisation underlying the interaction matrix
is not valid at the outset, so the runs exercise the convergence behaviour of
the schemes and not merely their behaviour near the solution. The iteration
cap is $600$ in all runs.

% ======================================================================

\subsection{Metrics}
\label{sec:metrics}
%%% (D) added; 4.6.2 refers to this subsection.

For each run the following are recorded: the evolution of the feature-error
norm $\norm{\mathbf{e}}^{2}$ against iteration; the final translational and
rotational pose errors; the adaptive gain and damping at convergence; and
whether the run converged, a run being taken as converged once the commanded
velocity falls below $10^{-3}\,\mathrm{m\,s^{-1}}$, for the reason given in
Section~\ref{sec:ablation-conditions}.
%%% (A) The failure definition is replaced by the criterion the chapter
%%% actually uses. (B) "the number of iterations to convergence" is
%%% removed: it is never reported in 4.7, and the sentence that excused it
%%% as secondary goes with it.

%%% (C) The standing of the feature error, declared where the metrics are
%%% declared. This paragraph is the reason to keep the subsection.
The final pose error is the primary criterion, and a word is needed on why it
is the only one. The feature error is the quantity the control law minimises,
and it is therefore a record of what the optimisation achieved rather than an
independent measure of success: under occlusion it is precisely the quantity
the corrupted measurements bias. Section~\ref{sec:occlusion} reports the
unweighted scheme attaining the \emph{smallest} final feature error of the
three variants while its pose error is larger by some two orders of
magnitude, the corrupted measurements having displaced the minimum of the
cost rather than raised its value there. The feature error is accordingly
reported throughout as evidence that a run converged --- that it reached a
minimum and settled --- and never as evidence of which minimum it reached.
Only the pose error, which is available in simulation and is independent of
the cost being minimised, can serve the latter purpose. This is a general
property of direct visual servoing and not an artefact of the present
configuration: any criterion computed from the same corrupted measurements
that drive the control inherits their bias.

In addition, the weight maps are recorded, that is, the weights $w_{j}$
displayed at their image positions, together with the structure measure $G$
displayed in the same way. These provide a direct qualitative check that the
weighting is suppressing the intended measurements and that the measure is
responding to structure rather than to intensity, and are reported alongside
the quantitative results in Section~\ref{sec:weight-maps}.
%%% Two changes. The structure maps are recorded too and are half of what
%%% Figure 4.1 and 4.7.3 show, so the list should name them. And the
%%% cross-reference is to 4.7.3 specifically, not to the whole of 4.7.

% ======================================================================
%  LABEL CLASH TO RESOLVE -- one line, but it will print "??" if missed.
%  Your text here refers to Section~\ref{sec:ablation-results}. My
%  correction to 4.5 refers to the same section as sec:results, and the
%  corrections to 4.4.3, 4.4.4 and 4.6.2 refer to its subsections as
%  sec:nominal (4.7.1), sec:occlusion (4.7.2) and sec:weight-maps (4.7.3).
%  Decide which label 4.7 actually carries and make the two agree; I have
%  used the subsection labels above because the references are to specific
%  results, not to the results section as a whole.
%
%  ONE THING I DID NOT REMOVE
%  "the adaptive gain and damping at convergence" is my addition, not
%  yours. 4.7.1 reports mu falling to order 1e-18 and the corrected 4.5
%  now rests its stability argument on that measurement. A quantity a
%  later section leans on should be declared as recorded here.
%
%  IF YOU WANT ONE MORE METRIC, the cheapest with the largest return is
%  the fraction of measurements carrying a weight below a quarter, per
%  iteration. You already compute it for Table 4.3 at iterations 1 and
%  600; logging it every iteration would turn the "rejection is not
%  immediate" claim of 4.4.3 from an inference drawn from two samples
%  into a curve. It is one accumulator in the loop.
% ======================================================================
